%%%%%%%%%%%%%%%%%%%%%%%%%%%%%%%%%%%%%%%%%%%%%%%%%%%%%%%%%%%%%%%%%%
% Pacotes e comandos acrescentados ao modelo tese-FT para este
% trabalho. O arquivo tese-FT.cls não foi alterado; tudo o que
% muda em relação ao modelo está aqui.
%%%%%%%%%%%%%%%%%%%%%%%%%%%%%%%%%%%%%%%%%%%%%%%%%%%%%%%%%%%%%%%%%%

% Tabelas que podem atravessar páginas, com cabeçalho repetido
\usepackage{xltabular}
% Listas com rótulo a), b), ... (objetivos específicos)
\usepackage{enumitem}
% Legendas acima de figuras e tabelas, como no texto original
\usepackage{caption}
\captionsetup{position=top,skip=6pt}
\setlength{\LTcapwidth}{\textwidth}

% Numeração contínua de figuras, tabelas e equações (Figura 1 a 12,
% Tabela 1 a 26, Equação 1 a 52), como no DOCX. Para numerar por
% capítulo (Figura 3.1), basta remover as três linhas a seguir.
\counterwithout{figure}{chapter}
\counterwithout{table}{chapter}
\counterwithout{equation}{chapter}

% Nos capítulos 3, 4 e 5 as seções são uma por dispositivo: 3A, 3B,
% com subseções 3A.1, 3A.2 etc. O comando é chamado logo depois do
% \chapter desses capítulos.
\newcommand{\secoesPorDispositivo}{%
  \renewcommand{\thesection}{\thechapter\Alph{section}}}

% Fonte e nota abaixo de figuras e tabelas
\newcommand{\notaabnt}[2]{%
  \par\nopagebreak{\footnotesize\setlength{\leftskip}{0pt}\setlength{\rightskip}{0pt}%
  \setlength{\parfillskip}{0pt plus 1fil}\noindent #1: #2\par}}
\newcommand{\fonte}[1]{\notaabnt{Fonte}{#1}}
\newcommand{\nota}[1]{\notaabnt{Nota}{#1}}

% Tabelas: espaçamento simples e corpo menor; nomes de arquivo longos
% podem quebrar depois do sublinhado. As tabelas curtas são flutuantes
% (table + tabularx); as longas usam o ambiente tabelaabnt, com
% xltabular, que atravessa páginas e repete o cabeçalho.
\newcommand{\estilotabela}[1]{%
  \setstretch{1}#1\renewcommand{\_}{\textunderscore\allowbreak}}
\newenvironment{tabelaabnt}[1]{%
  \par\begingroup\estilotabela{#1}%
  \setlength{\LTpre}{8pt}\setlength{\LTpost}{2pt}}%
  {\par\endgroup\vspace{8pt}}

% Espaço uniforme depois de ponto, como no texto original
% ("6 jun. 2018" sem espaço maior depois da abreviatura)
\frenchspacing

% Caracteres do texto corrido que o pdflatex não conhece
\DeclareUnicodeCharacter{03B1}{\ensuremath{\alpha}}
\DeclareUnicodeCharacter{03B2}{\ensuremath{\beta}}
\DeclareUnicodeCharacter{03B3}{\ensuremath{\gamma}}
\DeclareUnicodeCharacter{03B4}{\ensuremath{\delta}}
\DeclareUnicodeCharacter{03B5}{\ensuremath{\varepsilon}}
\DeclareUnicodeCharacter{03B7}{\ensuremath{\eta}}
\DeclareUnicodeCharacter{03B8}{\ensuremath{\theta}}
\DeclareUnicodeCharacter{03BA}{\ensuremath{\kappa}}
\DeclareUnicodeCharacter{03BB}{\ensuremath{\lambda}}
\DeclareUnicodeCharacter{03BC}{\ensuremath{\mu}}
\DeclareUnicodeCharacter{03C0}{\ensuremath{\pi}}
\DeclareUnicodeCharacter{03C1}{\ensuremath{\rho}}
\DeclareUnicodeCharacter{03C3}{\ensuremath{\sigma}}
\DeclareUnicodeCharacter{03C4}{\ensuremath{\tau}}
\DeclareUnicodeCharacter{03C6}{\ensuremath{\varphi}}
\DeclareUnicodeCharacter{03C8}{\ensuremath{\psi}}
\DeclareUnicodeCharacter{2113}{\ensuremath{\ell}}
\DeclareUnicodeCharacter{2192}{\ensuremath{\rightarrow}}
\DeclareUnicodeCharacter{2202}{\ensuremath{\partial}}
\DeclareUnicodeCharacter{2212}{\ensuremath{-}}
\DeclareUnicodeCharacter{221A}{\ensuremath{\surd}}
\DeclareUnicodeCharacter{2248}{\ensuremath{\approx}}
\DeclareUnicodeCharacter{2264}{\ensuremath{\leq}}
\DeclareUnicodeCharacter{2265}{\ensuremath{\geq}}
\DeclareUnicodeCharacter{22A5}{\ensuremath{\perp}}

% Grau de TCC de graduação. O modelo traz apenas \bsi e \tads, com
% os nomes dos cursos da FT; este comando segue o \bsi, com o nome
% do curso como argumento e sem a página da ficha catalográfica.
\makeatletter
\gdef\tccGraduacao#1{%
  \gdef\@textoaviso{\relax}%
  \if\@autora\relax
    \gdef\@degnamept{Bacharel em #1}%
    \gdef\@artigo{o}%
  \else
    \gdef\@degnamept{Bacharela em #1}%
    \gdef\@artigo{a}%
  \fi
  \gdef\@degnameen{Bachelor in #1}%
  \gdef\@degnamees{Licenciatura en #1}%
  \gdef\@monopt{Monografia}%
  \gdef\@monoes{Monografía}%
  \gdef\@monoen{Monograph}%
  \gdef\@nivelEnsino{Graduação}%
}
\makeatother

% Bibliografia: ajustes para reproduzir o formato da lista de
% referências do texto original
% uniquename=false: sem iniciais para distinguir autores de mesmo
% sobrenome (Kim, Wang, Ferguson), como nas citações do texto original
\ExecuteBibliographyOptions{useprefix=true,extradate=true,uniquename=false}
% intervalo de páginas com hífen (p. 79-96)
\DefineBibliographyExtras{brazil}{\protected\def\bibrangedash{-}}
% "Disponível em: URL", sem os sinais < >
\DeclareFieldFormat{url}{\bibstring{urlfrom}\addcolon\addspace\url{#1}}
% "p." minúsculo depois de ponto, em capítulos e anais
\DeclareFieldFormat[inproceedings,incollection]{pages}{p\adddot\addnbspace#1}
% "(ed.)" minúsculo: o estilo abnt põe o papel em maiúscula (\bibsentence)
\DeclareFieldFormat{editortype}{\printtext[parens]{#1}}
% nas patentes, a data só alimenta a citação; depósito e concessão
% estão no campo note
\AtEveryBibitem{\ifentrytype{patent}{\clearfield{year}\clearfield{month}\clearfield{day}}{}}

% Citação que combina referências da lista com trabalhos citados no
% texto que não constam da lista; tcc/gerar_latex.py a emprega quando
% isso acontece (hoje, todas as citações têm referência):
% \citecomsemref{Autor et al., 2017}{chave1,chave2}
\newcommand{\citecomsemref}[2]{%
  \AtNextCite{\renewcommand*{\postnotedelim}{\addsemicolon\space}}%
  \cite[#1]{#2}}
